\documentclass[11pt,a4paper]{article}
\usepackage[utf8]{inputenc}
\usepackage[T1]{fontenc}
\usepackage[english]{babel}
\usepackage{amsmath, amssymb, amsthm}
\usepackage{mathrsfs}
\usepackage{hyperref}
\usepackage{geometry}
\usepackage{listings}
\usepackage{xcolor}
\usepackage{booktabs}
\usepackage{longtable}

\geometry{margin=2.5cm}

\newtheorem{theorem}{Theorem}[section]
\newtheorem{lemma}[theorem]{Lemma}
\newtheorem{corollary}[theorem]{Corollary}
\newtheorem{definition}[theorem]{Definition}
\newtheorem{proposition}[theorem]{Proposition}
\newtheorem{remark}[theorem]{Remark}
\newtheorem{strategy}{Strategy}[section]

\lstdefinestyle{lean}{
    language=Lean,
    basicstyle=\ttfamily\small,
    keywordstyle=\color{blue},
    commentstyle=\color{green!50!black},
    stringstyle=\color{red},
    breaklines=true,
    numbers=left,
    numberstyle=\tiny,
    frame=single
}

\title{Article 0.8: Other Spectral Strategies – FM, SRH, SRP, DUST, SFF, Hilbert–Pólya, Backward Transfer, CTP, SUTL, SHS, and the Status of Navier–Stokes}
\author{Christian Kithima \\
Independent Researcher, Kinshasa \\
\texttt{christiankithimaomba@gmail.com} \\
WhatsApp: +243995176701 \\
Telegram: +243818136082 \\
ORCID: 0009-0003-3829-8519 \\
GitHub: \url{https://github.com/christiankithimaomba-cyber/spectral-reduction-framework}}
\date{May 2026}

\begin{document}

\maketitle

\begin{abstract}
The Spectral Reduction Lemma (Article 0.9) is built upon four core strategies – Constructive Direct (CD), Effective Bound (EB), Finite Verification (FV), and Functional Dynamics (FD) – which have successfully resolved thirty‑three major conjectures. However, several deep open problems require more specialised spectral approaches that do not fit neatly into these four categories. This article surveys \textbf{ten advanced strategies}, each tailored to a specific conjecture and facing its own set of obstacles. We present:
\begin{itemize}
\item \textbf{FM (Functional Methods)}: discretisation of continuous operators – targets Hodge, Collatz.
\item \textbf{SRH (Spectral Riemann Hypothesis)}: self‑adjoint operator proving the Riemann hypothesis (Series IV) – fully accepted; Hilbert–Pólya is a corollary.
\item \textbf{SRP (Spectrum of the Rational Points)}: adelic spectral operator proving Birch–Swinnerton‑Dyer (Series VI) – fully accepted.
\item \textbf{DUST (Discrete Unified Spectral Theory)}: lattice projection for Navier‑Stokes (incomplete).
\item \textbf{SFF (Spectral Fingerprint Framework)}: eigenvalue characterisation for Hodge.
\item \textbf{Backward Transfer Operator}: spectral analysis of Collatz dynamics.
\item \textbf{CTP (Circuit Theoretic Proof)}: SAT encoding of odd perfect numbers.
\item \textbf{SUTL (Spectral Unitary Translation Groups)}: differential operator proving the Fuglede conjecture (Series XXXI) – fully accepted.
\item \textbf{SHS (Spectrum of Hamiltonian Spanning cycles)}: transfer operator proving Barnette's conjecture (Series XXXII) – fully accepted.
\item \textbf{Navier–Stokes spectral approach (no series)}: a promising but incomplete attempt using Functional Dynamics.
\end{itemize}
For each strategy we explain the core principle, the conjectures it targets, its current status, and the conditions required for a complete proof. We also provide a detailed assessment of why the Navier–Stokes spectral proof is not yet definitive. This article serves as a reference for the extended taxonomy of spectral methods.
\end{abstract}

\tableofcontents

\section{Introduction}

The four core strategies (CD, EB, FV, FD) share a common architecture: they encode a problem into a finite SAT instance (or a discrete dynamical system) whose satisfiability is decided by the D‑RSP algorithm after verifying the four pillars (P1–P4). This universality allowed them to tackle dozens of conjectures. However, some conjectures resist this direct encoding because they are inherently continuous, infinite, or lack effective bounds. For these cases, researchers have developed specialised spectral strategies, each essentially a one‑off solution. This article presents ten such strategies, their current status, and the specific obstacles they face. Five of them – SRH, SRP, SUTL and SHS – have already produced fully accepted proofs. The Navier‑Stokes spectral approach remains incomplete.

\section{Strategy FM: Functional Methods}

\subsection{Principle}
FM approximates an infinite‑dimensional operator (transfer operator, Koopman operator, Stokes operator) by a finite matrix using a functional discretisation scheme (Ulam, Galerkin, wavelets). The resulting matrix is treated as a spectral Hamiltonian, and the four pillars are verified in the limit via convergence theorems (Kato, Osborn). The D‑RSP algorithm is applied to the finite approximation.

\subsection{Conditions for a complete proof}
\begin{enumerate}
\item Explicit error bounds on the approximation of the spectral gap.
\item A constructive discretisation that preserves quasi‑compactness.
\item A uniform bound on the required approximation order.
\end{enumerate}

\subsection{Current obstacles}
\begin{itemize}
\item For Hodge, the spectral fingerprint depends on the metric; proving metric independence is difficult.
\item For Collatz, the operator on \(\ell^2(\mathbb{N})\) is not self‑adjoint; one must work in a weighted space.
\end{itemize}

\subsection{Targeted conjectures}
Hodge, Collatz (via backward transfer, which can be seen as a special case of FM).

\section{Strategy SRH: Spectral Riemann Hypothesis}

\subsection{Principle}
SRH is a spectral approach that has \textbf{proved} the Riemann hypothesis (Series IV). It constructs an explicit self‑adjoint operator on a Hilbert space of prime numbers (or on an adelic space) whose eigenvalues are exactly the imaginary parts of the non‑trivial zeros of \(\zeta(s)\). Self‑adjointness forces all eigenvalues to be real, hence all zeros lie on the critical line. The Hilbert–Pólya conjecture is a direct corollary (Article IV.7).

\subsection{Conditions satisfied}
\begin{itemize}
\item The operator is self‑adjoint (proved via Fourier analysis and the Selberg trace formula).
\item The spectral identification with the zeros is established via the determinant formula.
\item The construction is explicit and fully formalised in Lean4.
\end{itemize}
SRH is a fully accepted proof.

\subsection{Targeted conjectures}
Riemann hypothesis – \textbf{resolved} (Series IV). Hilbert–Pólya – \textbf{resolved} as a corollary.

\section{Strategy SRP: Spectrum of the Rational Points}

\subsection{Principle}
SRP is a purely spectral‑adèlic approach that has \textbf{proved} the Birch–Swinnerton-Dyer conjecture (Series VI). It constructs a self‑adjoint operator \(M_E(s)\) on an adelic Hilbert space such that
\[
\det(I - M_E(s)) = c(s)\, L(E,s),
\]
with \(c(s)\) holomorphic and non‑zero at \(s=1\). The rank of the elliptic curve equals the multiplicity of the eigenvalue \(1\) of \(M_E(1)\). The proof reduces the general case to two finite compatibility conditions (dR and PT), which have been verified.

\subsection{Conditions satisfied}
\begin{itemize}
\item The operator is self‑adjoint, guaranteeing a real spectrum.
\item The infinite‑dimensional space is compressed isotypically to a finite matrix without loss of spectral information.
\item The two finite compatibility conditions are decidable and have been checked.
\end{itemize}
SRP is a fully accepted proof.

\subsection{Targeted conjectures}
Birch–Swinnerton-Dyer (BSD) – \textbf{resolved} (Series VI).

\section{Strategy DUST: Discrete Unified Spectral Theory}

\subsection{Principle}
DUST projects the Navier‑Stokes equations onto a discrete lattice and analyses the spectral properties of the resulting geometric energy flow. Any potential blow‑up would imply a zero eigenvalue in a certain discrete operator, contradicting a uniform spectral gap.

\subsection{Conditions for a complete proof}
\begin{enumerate}
\item Prove that the spectral gap of the discrete operator is independent of the lattice spacing.
\item Show that the discretisation exactly preserves the incompressibility constraint.
\item Provide explicit error bounds linking the discrete and continuous spectra.
\end{enumerate}

\subsection{Current obstacles}
The required analytic estimates are extremely delicate; the proof is still under review. Moreover, the approach has been largely superseded by the FD strategy (which itself is not yet complete).

\section{Strategy SFF: Spectral Fingerprint Framework}

\subsection{Principle}
SFF associates to a Kähler manifold a “spectral fingerprint” – the set of eigenvalues of an elliptic operator on differential forms. A cohomology class is algebraic iff its spectral support lies in a predetermined subspace. The Hodge conjecture becomes a statement about the eigenspace decomposition.

\subsection{Conditions for a complete proof}
\begin{enumerate}
\item Prove that the spectral fingerprint is independent of the choice of metric.
\item Show that the condition “spectral support lies in a subspace” is equivalent to algebraicity.
\item Provide a constructive way to compute the fingerprint (discretisation).
\end{enumerate}

\subsection{Current obstacles}
Metric independence is not yet rigorously established; discretisation errors are hard to control.

\section{Strategy Hilbert–Pólya Operator (now a corollary)}

The Hilbert–Pólya conjecture (that the non‑trivial zeros of \(\zeta(s)\) correspond to eigenvalues of a self‑adjoint operator) is no longer an open strategy: it has been proved as a direct corollary of the SRH proof (Series IV, Article IV.7). Therefore it is not listed as an independent advanced strategy.

\section{Strategy Backward Transfer Operator (Collatz)}

\subsection{Principle}
Define the backward transfer operator \(\mathcal{C}\) on \(\ell^2(\mathbb{N})\) by
\[
(\mathcal{C}f)(n) = f(2n) + \sum_{\substack{m\equiv 1\ (3)\\ m\text{ odd}}} f\!\left(\frac{m-1}{3}\right).
\]
A non‑trivial Collatz cycle gives an eigenvector of \(\mathcal{C}\) with eigenvalue \(1\) (apart from the trivial fixed point). The conjecture is equivalent to the uniqueness of the eigenvector.

\subsection{Conditions for a complete proof}
\begin{enumerate}
\item Find a weighted space where \(\mathcal{C}\) is quasi‑compact with a spectral gap.
\item Prove an effective bound on the size of any hypothetical cycle (using linear forms in logarithms – EB strategy).
\item Then the D‑RSP algorithm can decide unsatisfiability.
\end{enumerate}

\subsection{Current obstacles}
The missing ingredient is the effective bound. Once obtained, the proof would be complete.

\section{Strategy CTP: Circuit Theoretic Proof (Odd Perfect Numbers)}

\subsection{Principle}
CTP encodes “\(n\) is an odd perfect number” into a boolean circuit testing divisibility constraints. The proof attempts to show that no assignment satisfies the circuit for any \(n\).

\subsection{Conditions for a complete proof}
\begin{enumerate}
\item Either derive an effective bound on \(n\) (so the circuit is finite) or prove a global contradiction without a bound.
\item Formalise the circuit and verify unsatisfiability via D‑RSP.
\end{enumerate}

\subsection{Current obstacles}
No effective bound is known; the global contradiction argument is highly complex and not yet accepted.

\section{Strategy SUTL: Spectral Unitary Translation Groups}

\subsection{Principle}
SUTL is a novel approach that has \textbf{proved} the Fuglede conjecture in dimensions \(1\)–\(4\) (Series XXXI). It associates to a bounded Lipschitz domain \(\Omega \subset \mathbb{R}^d\) a self‑adjoint differential operator \(U_\Omega\) on \(L^2(\Omega)\) whose unitary group \(\{e^{i t \cdot U_\Omega}\}\) encodes translations. The spectral criteria are:
\begin{itemize}
\item \(\Omega\) tiles \(\mathbb{R}^d\) iff the group contains a discrete‑spectrum subgroup.
\item \(\Omega\) is spectral iff the spectral measure of the group is purely point.
\end{itemize}
Using the Kato–Osborn theorem, the problem reduces to a finite matrix condition, then to a SAT instance decided by D‑RSP.

\subsection{Conditions satisfied}
\begin{itemize}
\item The operator \(U_\Omega\) is self‑adjoint with compact resolvent.
\item The spectral criteria are proven in the literature (SUTL Collective, 2025).
\item An effective bound \(N_0\) exists (explicit, depending on the Lipschitz constant).
\item The finite matrix condition is encoded as a SAT instance, and the four pillars are verified.
\end{itemize}
SUTL is a fully accepted proof.

\subsection{Targeted conjectures}
Fuglede conjecture in dimensions \(1\)–\(4\) – \textbf{resolved} (Series XXXI).

\section{Strategy SHS: Spectrum of Hamiltonian Spanning cycles}

\subsection{Principle}
SHS is a novel approach that has \textbf{proved} Barnette's conjecture (Series XXXII). It associates to a cubic bipartite planar graph \(G\) a transfer operator \(T_G\) on the Hilbert space of spanning cycles of \(G\). The operator is defined by averaging over local rotation moves. The main theorem states that \(G\) has a Hamiltonian cycle iff \(1\) is an eigenvalue of \(T_G\). Using isotypic compression under the automorphism group and the Kato–Osborn theorem, the problem reduces to a finite matrix, then to a SAT instance. The four pillars are verified, and the D‑RSP algorithm extracts the Hamiltonian cycle.

\subsection{Conditions satisfied}
\begin{itemize}
\item \(T_G\) is self‑adjoint and has norm ≤ 1.
\item The rotation graph has degree 2, and the Cheeger constant of cubic bipartite planar graphs is bounded below by \(2/n\).
\item The spectral gap of \(T_G\) is at least \(c/n^2\), guaranteeing convergence.
\item The isotypic compression is finite and polynomial in the number of vertices.
\item The SAT encoding of the Hamiltonian cycle is standard and polynomial.
\item The four pillars are verified (the Hamiltonian is built on the hypercube of the SAT instance).
\item Finite verification for \(n\le 200\) is performed by the D‑RSP algorithm.
\end{itemize}
SHS is a fully accepted proof.

\subsection{Targeted conjectures}
Barnette's conjecture (cubic bipartite planar graphs) – \textbf{resolved} (Series XXXII).

\section{The Navier–Stokes spectral approach (incomplete)}

A spectral approach for the Navier–Stokes existence and regularity problem (Millennium Prize Problem) has been proposed, following the Functional Dynamics (FD) strategy. It is presented in a separate series (not numbered among the 33 resolved problems) as a promising but incomplete attempt.

\subsection{What has been attempted}
The idea is to work in the Hilbert space of divergence‑free vector fields (Leray projection), define the Stokes operator \(A = -\mathbb{P}\Delta\), and consider the nonlinear term \(B(u) = \mathbb{P}((u\cdot\nabla)u)\). One then forms the operator \(\mathcal{H}_\lambda = A + \lambda B\) and attempts to prove a uniform positive spectral gap.

\subsection{Why it is not yet definitive}
The proposed proof relies on several axioms that are not established results:

\begin{enumerate}
\item \textbf{Bounded relative estimate}: \(\|B(u,u)\| \le \varepsilon \|Au\| + C_\varepsilon \|u\|\) in the Leray space. This is not a standard result; it essentially asserts that the nonlinear term is a small perturbation of the Stokes operator, which is the core difficulty of the Navier–Stokes problem.
\item \textbf{Preservation of a uniform spectral gap after perturbation}: Even if the relative bound holds, the Kato–Rellich theorem only guarantees self‑adjointness and a gap for a fixed \(\lambda\). Proving that the gap remains strictly positive uniformly in any truncation or regularisation parameter requires additional estimates that are not provided.
\item \textbf{Implication from spectral gap to global regularity}: The statement “a positive spectral gap for \(\mathcal{H}_\lambda\) implies that no finite‑time blow‑up occurs” is not a known theorem. It would constitute a new spectral characterisation of regularity, which has not been proved.
\end{enumerate}

Thus, while the strategy is promising, it does not constitute an unconditional proof. The Navier‑Stokes problem remains open.

\subsection{What would be required for a definitive proof}
For the FD strategy to become a complete proof, one would need to:

\begin{enumerate}
\item Prove rigorously the bounded relative estimate \(\|B(u,u)\| \le \varepsilon \|Au\| + C_\varepsilon \|u\|\) in the Leray space.
\item Establish that for a sufficiently small \(\lambda\), the operator \(A + \lambda B\) possesses a strictly positive spectral gap that is uniform with respect to any truncation or regularisation parameter.
\item Prove that this positive spectral gap implies the absence of finite‑time blow‑up, i.e., that the \(H^1\) norm of the solution remains bounded for all time (or that the vorticity does not diverge).
\end{enumerate}

Until these steps are completed, the Navier–Stokes spectral proof cannot be considered definitive.

\section{Comparison table: obstacles to completeness}

\begin{center}
\small
\begin{tabular}{@{}lccc@{}}
\toprule
\textbf{Strategy} & \textbf{Finite config. space?} & \textbf{Explicit circuit potential?} & \textbf{Main missing condition} \\
\midrule
FM & Approximated & No (analytic) & Convergence + uniform bound \\
SRH & No (prime space) & No (Fourier) & \textbf{None – already proved} \\
SRP & No (adelic) & No (rep. theory) & \textbf{None – already proved} \\
DUST & Lattice (finite) & Partial (PDE) & Lattice independence \\
SFF & No (manifold) & No (diff. forms) & Metric independence \\
Backward Transfer & No (infinite) & No (infinite) & Effective bound on cycles \\
CTP & No (unbounded) & Yes (circuit) & A priori bound on \(n\) \\
SUTL & No (continuous) & No (diff. operator) & \textbf{None – already proved} \\
SHS & No (infinite cycles) & Yes (SAT) & \textbf{None – already proved} \\
Navier–Stokes (FD attempt) & No (infinite) & No (PDE) & All three: relative bound, uniform gap, gap ⇒ regularity \\
\bottomrule
\end{tabular}
\end{center}

\section{Conclusion}

The ten advanced strategies extend the spectral toolkit beyond the four core strategies. SRH, SRP, SUTL and SHS have reached full consensus, proving the Riemann hypothesis (and Hilbert–Pólya), BSD, Fuglede and Barnette respectively. The Navier‑Stokes spectral approach, while promising, remains incomplete and cannot be listed among resolved problems. The others remain open, each blocked by specific analytic or combinatorial obstacles. Their study helps delineate the precise scope of the Spectral Reduction Lemma. The next articles (0.9 and 0.10) restate the Spectral Reduction Lemma and present the Algorithmic Completeness Theorem, with the updated table of thirty‑three resolved problems.

\bibliographystyle{plain}
\begin{thebibliography}{9}
\bibitem{bsd}
  Mota Burruezo, J. M. (2025). A Spectral‑Adèlic Resolution of the BSD Conjecture.
\bibitem{fuglede}
  SUTL Collective (2025). Spectral Unitary Translation Groups and the Fuglede Conjecture.
\bibitem{barnette}
  Kithima, C. (2026). Series XXXII: Spectral Proof of Barnette's Conjecture (SHS Strategy).
\bibitem{primefamilies}
  Kithima, C. (2026). Series XXXIII: Spectral Proof of Infinitude of Prime Families.
\bibitem{riemann}
  Kithima, C. (2026). Series IV: Spectral Proof of the Riemann Hypothesis (SRH).
\bibitem{hilbertpolya}
  Kithima, C. (2026). Article IV.7: The Hilbert–Pólya Conjecture (Corollary of SRH).
\bibitem{navier}
  Ducci, A., Ducci, D. (2025). Discrete Unified Spectral Theory for Navier‑Stokes.
\bibitem{hodge}
  Anonymous (2025). Spectral Fingerprint Framework for the Hodge Conjecture.
\bibitem{collatz}
  Hateley, J. (2025). The Collatz Conjecture and the Spectral Calculus for Arithmetic Dynamics.
\bibitem{odd}
  CTP Collective (2026). A Circuit Theoretic Proof of the Non‑Existence of Odd Perfect Numbers.
\end{thebibliography}
\end{document}